%! TEX root = ../main.tex

\chapter{Discrimination of Photons and Protons}
\label{chap:chapter-4}

The attempt to identify the primary particle type from a (possibly incomplete) 
\EAS shower footprint detected by the \SD is a difficult task, as layed out in 
\cref{ssec:different-primaries}. While shower-to-shower fluctuations complicate 
predictions based on simplified cascade models, a general trend holds true when
comparing the \EASs initiated from different primaries; The muonic fraction of
the \EAS grows with the 

muonic fraction grows with log A

\section{Monte Carlo Datasets}

Confirmed photon \CRs, detected e.g.~by the \LHAASO have reconstructed primary 
energies of $<\SI{10}{\peta\eV}$ 
\cite{lhaasocollaborationUltrahighenergyRayBubble2024}, less than 
\SI{10}{\percent} of the lower bound on the energy considered in this work. No 
photon candidates have ever been measured in the search space studied here.

With this lack of experimental evidence we must base our analysis on simulations
of \EASs. The following sections detail how these simulations are performed, and
present the datasets of simulated air showers used throughout 
\cref{chap:chapter-4}.

The code used in generating the simulation datasets are available at 
Ref.~\cite{AugerobservatorySandBoxesPaulfilip2026}.

\subsection{Air Shower Simulation}

\texttt{CORSIKA v7.8010} \cite{engelNextGenerationCORSIKA2018} is used to 
simulate air showers for both photon and proton primaries. In either case, 
primary particles with a true energy $E_\mathrm{MC}$ are injected with a heading
$(\phi_\mathrm{MC},\,\theta_\mathrm{MC})$ into the upper atmosphere, and the 
resulting particle cascade is simulated using \texttt{EPOS\_LHC} 
(Ref.~\cref{sssec:hadronic-interaction-models}, \cite{pierogEPOSLHCTest2015}), 
with \texttt{FLUKA  2011.2c} \cite{theflukacollaborationFLUKACodeOverview2024} 
handling particle interactions below a center-of-mass energy of 
\SI{80}{\giga\eV}.

\subsubsection{Particle Distribution}

Particles are injected at an atmospheric depth of 
\SI{0}{\gram\per\centi\meter\squared}, with the parametrized air density profile
emulating typically measured atmospheric conditions in Malargüe in September, as
described in further detail in 
Ref.~\cite{abreuDescriptionAtmosphericConditions2012}. 

The heading of the primary particle is randomized with the distribution of 
$(\phi_\mathrm{MC},\,\theta_\mathrm{MC})$ mimicking a uniform flux over the sky 
as observed by a flat detector. Specifically,
$\phi_\mathrm{MC}\sim\mathcal{U}(\SI{-180}{\degree},\,\SI{180}{\degree})$, and 
$\sin(2\,\theta_\mathrm{MC})\sim\mathcal{U}(0,\,0.766)$. The upper bound on the 
latter parameter corresponds to a maximum simulated zenith angle of 
$\SI{65}{\degree}$. The distribution of simulated primary zenith angles is shown
in \cref{fig:zenith-mc-distribution}.

The primary energy of injected particles is 
$\log_{10}\left(E_\mathrm{MC}\,/\,\mathrm{eV}\right)\in[16.5, 19.0]$, covering a
broad range from full inefficiency,\footnote{Meaning no \EAS below this energy 
is ever observed in the \SD-750.} to an upper energy where the size of the array
becomes a bottleneck in detecting the faint \CR fluxes. For an individual 
\texttt{CORSIKA} run the energy range of the particle is fixed to one tenth of a
decade in logarithmic energy. Within this bin, the energy is randomized such 
that $E^\mathrm{bin}_\mathrm{MC}\sim E$, meaning that higher values are picked 
more likely than lower values. In order to obtain sufficient sample sizes across
all energy bins lower energy showers are simulated more frequently than higher 
energy showers. Both considerations influence the distribution of simulated 
primary energies, as shown in \cref{fig:energy-mc-distribution}.

\begin{figure}[t]
  \centering
  \subfloat[]{\includegraphics[height=0.36\textwidth]{discrimination/zenith_mc_distribution.png}
  \label{fig:zenith-mc-distribution}
  }\hspace{0.2cm}
  \subfloat[]{\includegraphics[height=0.36\textwidth]{discrimination/energy_mc_distribution.png}
  \label{fig:energy-mc-distribution}
  }
  \caption[]{\subref{fig:zenith-mc-distribution}  \subref{fig:energy-mc-distribution} }
  \label{fig:mc-distributions}
\end{figure}


\todo{simulate wider range of energies? Update plots}

In order to be computationally efficient the simulation framework employs 
\emph{thinning} \cite{saffiAnalysisSimulationReconstruction2020}, which assigns
weights to individual particles in the shower cascade. A particle with weight 
$W > 1$ consequently approximates $W$ copies of itself, reducing the number of 
tracked particles, and cutting down on processing time. The threshold for the 
thinning algorithm to group particles together is energy dependent. This 
threshold is set dynamically to an optimal value of $10^{-6}\,E_\mathrm{MC}$ 
\cite{billoirSamplingProcedureRegenerate2008}.

The simulation of the shower cascade halts at an atmospheric depth corresponding
to an altitude of $\SI{1400}{\meter}\,\mathrm{a.s.l.}$, the height of the \PAOs 
detectors. At this point, the properties of the remaining particles in the 
simulated shower front are handed to an \Offline simulation that generates (\SD)
detector signals from the particle distribution on ground. In order to further 
optimize computational usage, we repeat this process multiple times per shower, 
varying the position of the \SD detector relative to the \EAS every time. Each 
\texttt{CORSIKA} simulation is passed to \Offline at most $50$ times, or until 
the \EAS has been reconstructed successfully $n_\mathrm{success} = 10$ times. 

\subsubsection{Detector Signal}

For a single reconstruction run, the shower core, referring to the intersection 
between the shower axis (fixed by $\phi_\mathrm{MC},\,\theta_\mathrm{MC}$) and 
the detector plane is picked at random. An arbitrary \SD-750 station is chosen, 
and the shower core is picked as a random point from a unit rectangle centered 
around this station. In this fashion, the distribution of simulated shower cores
mimics the uniform\footnote{Non-uniformities are caused by detector 
inefficiencies, see e.g.~Ref.~\cref{TODO}} distribution of \EASs core 
distributions the surface of the Earth. 

Once the shower core is picked, the entire geometry of shower and detector is 
fixed. Particles (with weight $W$) that fall within a predefined volume around 
an \SD-station are \emph{dethinned} 
\cite{saffiAnalysisSimulationReconstruction2020}, and particle trajectories
through the \WCD, \SSD, and \UMD, as well as resulting material interactions are
simulated using \texttt{Geant4} \cite{allisonRecentDevelopmentsGeant42016, 
allisonGeant4DevelopmentsApplications2006, 
agostinelliGeant4aSimulationToolkit2003}. The trajectory of light created in the
scintillating material of the \SSD and \UMD, as well as in the Cherenkov cone of
relativistic muons in the \WCD is propagated to the different sensor components,
eventually resulting a simulated detector response to the virtual particle flux.

On top of this, electronic noise induced by the \DAQ readout chain, as well as 
low-energy atmospheric background muons are injected into the detector 
simulation to increase the degree of realism of the calculations.

\subsection{Event Reconstruction}

The simulated detector response from the \WCD, \SSD, and \UMD is reconstructed 
according to the standard reconstruction procedure defined in 
\cref{ssec:sd-reconstruction} and \cref{ssec:umd-reconstruction}. The module 
configuration settings used for the bulk of the analysis are the default values 
used in the \texttt{icrc2025-test7} release of \Offline 
\cite{Icrc2025test7Release2025}. Influences due to changes in the configuration 
files will be examined in \todo{hyperparameter optimization chapter}.

\subsection{Reconstruction Efficiency}

An important characteristic in the subsequent analysis is the reconstruction 
efficiency of the \SD(-750) array. Together with the distribution of a priori 
present \CR particles (i.e. the \CR spectrum), it determines the energy-, and 
zenith distribution of \CR observations. In particular, 
$P_\mathrm{rec}(E_\mathrm{MC}, \theta_\mathrm{MC})$\footnote{We assume 
$P_\mathrm{rec}(E_\mathrm{MC}, \theta_\mathrm{MC}, \phi_\mathrm{MC}) = 
P_\mathrm{rec}(E_\mathrm{MC}, \theta_\mathrm{MC})\,\,\,\forall\,\phi_\mathrm{MC}
\in [-180\,^\circ,\,180\,^\circ]$} gives the probability that a \CR event with 
true energy $E_\mathrm{MC}$ and inclination $\theta_\mathrm{MC}$ is successfully
reconstructed. This is to say that the \texttt{LDFFinderKG} 
(Ref.~\cref{sssec:reconstruction-ldf-finder}) runs without emitting errors, and 
the event can be assigned a reconstructed energy $E_\mathrm{rec}$. 

Importantly, this probability differs from the often quoted array trigger 
efficiency $P_\mathrm{T3}$ 
(e.g.~\cite[Fig.~7]{abrahamTriggerApertureSurface2010}). The detection of an 
event does not guarantee its successful reconstruction, on which the problem at 
hand, the identification of exotic events in reconstructed data, relies on. 
$P_\mathrm{rec}$ is thus the appropriate choice for normalizing different 
detector characteristics in this analysis.

We employ a parametrization of $P_\mathrm{T3}$/$P_\mathrm{rec}$ from \todo{cite}, namely

\begin{equation}
\label{eq:rec-efficiency-parametrization}
P_\mathrm{rec}(E_\mathrm{MC}, \theta_\mathrm{MC};\,\alpha(\theta_\mathrm{MC}), \beta) = \frac{1 + \erf\left(E_\mathrm{MC} - \alpha(\theta))\,/\,\beta\right)}{2}, 
\end{equation}

where

\begin{equation}
\label{eq:efficiency-theta-dependence}
\alpha(\theta_\mathrm{MC};\,a_0,\,a_1,\,a_2,\,a_3) = a_0 + a_1\,x + a_2\,x^2 + a_3\,x^3,\qquad x = \cos^2\,\theta_\mathrm{MC}\,.
\end{equation}

The shape parameters $a_i$ and $\beta$ are determined from an unbinned maximum 
likelihood optimization tailoring \cref{eq:rec-efficiency-parametrization} to 
the reconstruction efficiencies calculated during the simulation procedure. 
These follow (for a single simulation with energy $E_\mathrm{MC}$ and zenith 
$\theta_\mathrm{MC}$) by building the ratio of number of successful 
reconstruction $n_\mathrm{success}$ versus number of total attempted 
reconstructions. The results of the optimization, both for $p$ and $\gamma$ 
primaries is displayed in \todo{table}. The visual correspondence between the
simulated reconstruction efficiency and the parametrized reconstruction 
efficiency is shown in \todo{figure}.

\subsubsection{Full Efficiency Transition}



% important to not overestimate reconstruction efficiency
% cross check from full efficiency transition w/ arrival direction?

\subsection{Energy Assignment}
\subsubsection{Simulated Protons}
\subsubsection{Simulated Photons}
\subsection{Unified Energy Scale}

\section{Validation with Past Studies}

\subsection{Reconstruction Efficiency}

\section{μLDF-based Discrimination}
% trigger, zenith, station selection, muLDF dependence

\subsection{Parametric Dependence}
\subsection{Selection Dependence}
\subsection{Trigger Dependence}
\subsection{Zenith Dependence}

\section{Validation on Burnt Sample}

\section{Further Improvements}
\subsection{Template Fits}
\subsection{Graph Neural Networks}

\section{Summary}
